\section{Scheduler Design}\label{sec:method}
\subsection{Prepare shared state once per timestamp}
Policy masks are reusable while service requirements and node capabilities
remain unchanged. For each timestamp, the scheduler constructs node features
and batches model inference over the nodes required by the event batch.
The resulting risk and headroom vectors are reused by individual events.
Computing a risk score does not admit a node: policy remains a mandatory
condition on every selected source or destination.

The reported predictor uses 22 features covering current CPU and memory,
projected values, joint utilization, headroom, two preceding snapshots,
changes, and short-window statistics. These features use three consecutive
snapshots ending at $t$. The feature families describe the predictor's inputs but do not fully
specify its implementation. The serialized feature schema, individual
feature formulas, and model hyperparameters are needed for exact
reproduction and are not included with this draft.

\subsection{Shortlist before final screening}
For a source that fails admission, \cage orders $\mathcal B_e$
lexicographically by $(H(n),g_{n,t},\operatorname{id}(n))$ and takes its
first $k$ nodes. It then applies the risk and headroom screens to this
shortlist. If none pass, it screens the remaining base candidates.
Among accepted nodes it selects the minimum
\begin{equation}\label{eq:ranking}
\left(H(n),p_{n,t},g_{n,t},\operatorname{id}(n)\right).
\end{equation}
Thus, historical use is the primary ranking key, and predicted risk breaks
ties in that key. The first shortlist is drawn from the policy-and-batch
candidate set, not from $\Fset_e$. This distinction matters: shortlisting
before screening can choose a different node from ranking all feasible
nodes. Equality with full ranking is an empirical result for the reported
workloads, not a general property of the algorithm.

If all base candidates fail, the scheduler returns an infeasible decision.
Search expansion guarantees that a candidate passing the fixed snapshot
checks is not missed merely because it lies outside the initial shortlist.
It does not guarantee that every event has a feasible destination.

\subsection{Final check and state update}
Algorithm~\ref{alg:scheduler} consolidates the decision path. Before
recording a migration, the scheduler repeats the policy, risk, headroom,
active-source, and reservation checks. A failed recheck cannot fall through
to a committed decision. In the replay, telemetry is fixed within a batch
and decisions are processed sequentially. A live controller would also need
atomic resource reservation and a defined retry policy for changing state.

\begin{algorithm}[t]
\caption{Placement decisions for one timestamp}\label{alg:scheduler}
\KwIn{Ordered events $\Eset_t$; policy masks; node snapshots; history $H$}
Prepare $p_t,g_t$ once; form $A_t$; set $U_t\leftarrow0$\;
\ForEach{$e=(i,s,n_s,t)\in\Eset_t$}{
  $\pi(e)\leftarrow\varnothing$\;
  \If{$L(s,n_s)=1$ and $p_{n_s,t}\leq\theta$ and $g_{n_s,t}\leq1$}{
    $\pi(e)\leftarrow n_s$; $H(n_s)\leftarrow H(n_s)+1$\;
    \textbf{continue}\;
  }
  Form $\mathcal B_e$ using Eq.~\eqref{eq:base}\;
  $K\leftarrow$ first $k$ nodes by $(H,g,\operatorname{id})$\;
  $D\leftarrow\{n\in K:p_{n,t}\leq\theta,\ g_{n,t}\leq1\}$\;
  \If{$D=\varnothing$}{
    $D\leftarrow\{n\in\mathcal B_e\setminus K:p_{n,t}\leq\theta,\ g_{n,t}\leq1\}$\;
  }
  \If{$D=\varnothing$}{\textbf{continue}\;}
  $n^*\leftarrow$ minimum of $D$ by Eq.~\eqref{eq:ranking}\;
  Recheck all conditions in Eq.~\eqref{eq:feasible} for $n^*$\;
  \If{any recheck fails}{\textbf{continue}\;}
  $U_t(n^*)\leftarrow U_t(n^*)+1$\;
  $H(n^*)\leftarrow H(n^*)+1$; $\pi(e)\leftarrow n^*$\;
}
\Return{$\pi$}\;
\end{algorithm}

The decision procedure preserves its admission predicates under a fixed
snapshot: retention tests all three source conditions, and migration is
recorded only after every destination condition passes. This is a statement
about the algorithm's checks. It says nothing about the accuracy of the
predictor, the truth of supplied policy attributes, or future resource use.

\subsection{Cost}
For $R_t$ nodes requiring inference and a fixed trained model, feature
construction and inference take $O(R_t)$ work. Building a candidate set is
$O(M)$ per migrating event. A heap-based top-$k$ implementation can select
the shortlist in $O(M\log k)$, then rank accepted candidates in $O(k)$;
expansion adds an $O(M)$ scan. A full sort instead costs $O(M\log M)$.
These are implementation bounds; timing results alone do not establish
which selection primitive was used. State requires $O(M)$ space plus $O(PM)$ for $P$ policy masks.
Shortlisting reduces candidate-ranking work, not the number of nodes whose
features necessarily need to be constructed.
